\section{Transport and the graph variance bound}\label{sec:transport}

We now turn the quadratic property into a variance bound for the graph of
Section~\ref{sec:graph}. A signed flow compares adjacent conditional means;
decomposing variance over successive conditions gives a bound on the
transversals, and the repair map extends it to the full graph. All large
slots are summed with their balances throughout. The remaining slots are
exposed in row-major order. Fix a real-valued function $H$ on the feasible
small-state graph.
The flow-energy comparison is the standard quadratic electrical-network
method; see Aldous and Fill~\cite[Section~3.7.1]{AF2002}. The recursive
signed demands and their bound for integer slots are constructed here.
A \emph{context} $\sigma$ fixes balanced profiles on a prefix of small slots.
Write $\mathcal T_\sigma$ for the transversals with this context and
$z_\sigma=\sum_{X\in\mathcal T_\sigma}f(X)$.  If $s$ is the next slot,
its child $j$ fixes $(x_s,y_s)=(j,U-j)$, and has mass $z_j$ and weighted
mean $H_j$.  Means are used only for positive masses.

Let $E_\sigma(H)$ be the sum of
$\min(f(X),f(Y))(H(X)-H(Y))^2$ over unordered \emph{unit-exchange} edges
whose endpoints respect $\sigma$.  In particular $E_\sigma(H)\le E(H)$;
repair edges need not be included in $E_\sigma$.  A superscript $\eta$ on
masses or means, and $E_{\eta,\sigma}$ on this energy, will indicate the
softened weights from Definition~\ref{def:soft-weight}.  For the softened
energy we include all box-admissible transversal and defect patterns with
the context, including states whose limiting weight is zero.

\subsection{Child weights and adjacent means}

\begin{lemma}[Child log-concavity]\label{lem:child-logconcavity}
At every context with positive transversal mass, the child masses
$(z_j)_{j=0}^U$ are log-concave and their positive support is an integer
interval.  If $z_a,z_b>0$ and $a\le j\le b$, then
$z_j\ge\min(z_a,z_b)$.
\end{lemma}
\begin{proof}
For $0<\eta<1$ every child mass is positive.  For $1\le j<U$, display
the next pair as $(j+1,U-j+1)$ and sum every other unexposed slot with its
balance.  The display is in its box and has two excess units.  Its
two-removal matrix is
\[
 \begin{pmatrix}z_{j-1}^\eta&z_j^\eta\\
                 z_j^\eta&z_{j+1}^\eta\end{pmatrix}.
\]
Proposition~\ref{prop:signature} gives at most one positive eigenvalue.
The trace is positive, so the determinant is nonpositive:
$(z_j^\eta)^2\ge z_{j-1}^\eta z_{j+1}^\eta$.
Consequently the successive differences of $\log z_j^\eta$ are
nonincreasing, and for $a<j<b$,
\[
 z_j^\eta\ge
 (z_a^\eta)^{(b-j)/(b-a)}(z_b^\eta)^{(j-a)/(b-a)}.
\]
Finite sums converge as $\eta\downarrow0$.  The local inequalities give
log-concavity in the limit; the displayed interpolation gives positivity
between positive endpoints and the asserted lower bound.
\end{proof}

\begin{lemma}[Conditional transport]\label{lem:transport}
Let $\sigma$ be a row-major context, and let children $l-1,l$ of its
next slot have positive masses $z_-,z_+$ and means $H_-,H_+$.  Then
\begin{equation}\label{eq:conditional-transport}
 \min(z_-,z_+)(H_--H_+)^2
 \le A_0 E_\sigma(H),\qquad A_0=2+2d(U+1)^2.
\end{equation}
\end{lemma}
\begin{proof}
We first work at a fixed $\eta>0$. Extend $H$ by arbitrary fixed finite
values to the box-admissible graph patterns whose limiting weights are
zero, and suppress $\eta$ on masses. We seek a signed flow with divergence
$f_\eta(X)/z_-$ at each transversal in the minus child,
$-f_\eta(X)/z_+$ at each transversal in the plus child, and zero elsewhere.
For such a flow $J$, summation by parts and Cauchy--Schwarz give
\[
 (H_-^\eta-H_+^\eta)^2
 \le \left(\sum_{XY}\frac{J(X,Y)^2}
                    {\min(f_\eta(X),f_\eta(Y))}\right)
                         E_{\eta,\sigma}(H),
\]
where each unit-exchange edge is oriented once. Thus it is enough to
bound the flow's resistance energy. We build its signed demands by
successively exposing the remaining small slots, and charge its energy
to a quadratic potential along that recursion.

Display the special slot $s$ as $(l,U+1-l)$, where $1\le l\le U$.
Removing its first or second coordinate gives the minus or plus child,
respectively. The flow uses these endpoint transversals and defects
with $s$ positive and one other unexposed small slot negative.

\emph{Balanced recursion.}
A recursion node $\nu$ displays the special pair and some subsequent
small pairs, the latter all with occupancy $U$. Let $p$ be this display,
and call its coordinates the old indices. With the context included,
and all undisplayed slots assigned their balanced occupancies, the total
is $T+1$, where $T=\sum_s B_s$. For an old index $i$, let $s_i$ be the
total softened weight obtained by removing one unit there and summing
all undisplayed slots with their balances. Invalid removals have weight
zero. Assign coefficients $h_i$ satisfying
\begin{equation}\label{eq:balanced-demands}
 \sum_i s_i h_i=0,\qquad \Phi_\nu=\sum_i s_i h_i^2.
\end{equation}
The quantity $s_i h_i$ is the signed demand on that one-removal fiber;
the first equality requires total demand zero at the node. At the root
the two coefficients are $1/z_-$ and $-1/z_+$, giving the prescribed
endpoint demands and root potential $1/z_-+1/z_+$.

To process the next ordinary slot $v$, branch over its balanced profiles
$(j,U-j)$, $0\le j\le U$.  In the notation of
Lemma~\ref{lem:pair-summation}, use the auxiliary occupancy $D'=U+1$.
Then $S(j)$ is the vector of old-hole weights in branch $j$, and $u_j$
is the no-old-hole weight with pair $(j-1,U-j)$.  Thus $u_1,\ldots,u_U$
are positive and $u_0=u_{U+1}=0$.  Since $\sum_j S_i(j)=s_i$, define
\begin{equation}\label{eq:transport-g}
 g_j=\frac{\sum_{a=0}^{j-1}h^T S(a)}{u_j}\quad(1\le j\le U),
 \qquad g_0=g_{U+1}=0.
\end{equation}
Balance gives $h^TS(j)=u_{j+1}g_{j+1}-u_jg_j$, including both endpoints.
In child $j$, retain all old coefficients and give the two new
coordinates coefficients $g_j,-g_{j+1}$.  Their one-hole masses are
$u_j,u_{j+1}$, so the child is balanced.  Summing child potentials
conserves the old terms and adds exactly
\begin{equation}\label{eq:potential-increment}
 \sum_{j=0}^U\Phi_{\nu j}-\Phi_\nu
 =2\sum_{j=1}^Uu_jg_j^2
 \le -h^TV_{\nu,v}h.
\end{equation}
Here $V_{\nu,v}$ is the matrix of two-old-hole weights with $v$ displayed
at $(a,U+1-a)$, summed over $1\le a\le U$ with multiplier
$a(U+1-a)$ and over all other undisplayed slots with their balances.
The inequality is precisely Equation~\eqref{eq:pair-weighted}.
Each individual display before the two removals has total $T+2$.

\emph{A future quadratic form.}
For any slot $t$ still to be processed, define
$Q_t(\nu)=h^TV_{\nu,t}h$ by the same rule.  If a different slot $v$ is
split now, we claim
\begin{equation}\label{eq:future-monotonicity}
 \sum_{j=0}^U Q_t(\nu j)\ge Q_t(\nu).
\end{equation}
For an explicit common summation kernel, put
\[
 K_t(r,w)=\sum_{a=1}^U a(U+1-a)
       \sum_{b}F_\eta\bigl(\sigma,r,w,(a,U+1-a),b\bigr).
\]
Here $r$ assigns the old coordinates, $w$ assigns the current pair $v$,
and $b$ ranges over the Cartesian product of balanced profiles of every
remaining slot other than $t$, including all large slots.
As usual, an invalid coordinate assignment contributes zero.  This is
a single fixed summation domain; all penalties are evaluated on the
full configurations.  Every kernel entry used below has total $T$.

The old--old terms in the children partition the parent's balanced sum
over $v$.  For a mixed term, removing old index $i$ and the first new
coordinate in branch $j$ gives
$K_t(p-e_i,(j-1,U-j))$ with coefficient $2h_i g_j$.
Removing $i$ and the second new coordinate in branch $j-1$ gives the
identical kernel entry with coefficient $-2h_i g_j$.
These terms cancel.  All legal resulting pairs have occupancy $U-1$;
the unmatched formal endpoint removals are invalid and have weight zero.
For the new--new terms, each resulting pair $(k,U-2-k)$, $0\le k\le U-2$,
has exactly the following contributions:
\begin{center}
\begin{tabular}{lll}
\toprule
Removed coordinates & Branch & Coefficient\\
\midrule
First twice & $k+2$ & $g_{k+2}^2$\\
First and second & $k+1$ & $-2g_{k+1}g_{k+2}$\\
Second twice & $k$ & $g_{k+1}^2$\\
\bottomrule
\end{tabular}
\end{center}
They all multiply $K_t(p,(k,U-2-k))$.  Hence the exact identity is
\begin{equation}\label{eq:future-square}
 \sum_{j=0}^U Q_t(\nu j)-Q_t(\nu)
 =\sum_{k=0}^{U-2}K_t(p,(k,U-2-k))
                      (g_{k+2}-g_{k+1})^2\ge0.
\end{equation}
All omitted boundary removals are invalid.  Auxiliary entries with two
holes in an ordinary pair need not be graph states: they belong to the
full softened total-$T$ slice where the signature proposition applies.
No signature property of the weighted sum $V_{\nu,t}$ is being asserted;
Equation~\eqref{eq:potential-increment} and the explicit square identity
are the properties of that sum that are used.

\emph{The root bound.}
Iterating Equation~\eqref{eq:future-monotonicity} shows that the sum of
$Q_t$ at the level just before processing $t$ is at least its root value.
Thus Equation~\eqref{eq:potential-increment}, summed over levels, bounds
the total leaf potential by the root potential minus the sum of the
root $Q_t$'s.  Let $c_t=c_{st,l}^{\eta,\sigma}$ be the total weight of
defects respecting $\sigma$ with special pair $(l-1,U-l)$ negative,
slot $t$ positive, and every other small slot balanced.
These auxiliary defects differ from the special-positive defects used by
the leaf flow: the off-diagonal entry removes both special coordinates
from $(l,U+1-l)$, leaving $(l-1,U-l)$.
At the root, the off-diagonal entry of $V_{\nu,t}$ is this sum with
multiplier $a(U+1-a)$ on the positive pair $t$, so it is at most
$(U+1)^2c_t$.  Its diagonal entries are nonnegative.  Substituting the
root coefficients and dropping their nonpositive contributions to
$-Q_t$ gives
\begin{equation}\label{eq:leaf-potential}
 \sum_{\nu\text{ leaf}}\Phi_\nu
 \le \frac1{z_-}+\frac1{z_+}
       +\frac{2(U+1)^2}{z_-z_+}\sum_{t>s}c_t.
\end{equation}
The sum ranges only over the other unexposed small slots.  With no such
slots the recursion consists of its root and this bound still holds.

\emph{Leaf flows and their divergences.}
At a leaf let $D$ denote the entire small-coordinate display, including
the fixed context. For its unfixed coordinate indices $i$, the valid
one-removal states $X_i=D-e_i$ have
weights $s_i=f_\eta(X_i)>0$, and form a clique of unit exchanges.
Choose a maximum-weight vertex as hub.  On the edge from each nonhub
$X_i$ to the hub send signed flow $s_i h_i$; its conductance is exactly
$s_i$.  The flow's divergence at $X_i$ is $s_i h_i$, and balance gives
the same prescribed divergence at the hub.  Its resistance energy is
at most $\sum_i s_i h_i^2=\Phi_\nu$.
Different leaf cliques have disjoint edge sets: the coordinatewise
maximum of the endpoints recovers $D$, which uniquely determines the
leaf.  Their flows can therefore be added without an energy multiplier.

An endpoint transversal has a unique leaf representation.  Its ordinary
balanced profiles determine all branches, and its special hole is in
the first coordinate for the minus child and the second for the plus
child.  Its coefficient remains $1/z_-$ or $-1/z_+$ throughout.
Every other used state has special slot positive and a unique ordinary
negative pair, say $(a,U-1-a)$ in slot $v$.  It has exactly two leaf
representations: fill the first coordinate, giving branch $a+1$ and
coefficient $g_{a+1}$, or fill the second, giving branch $a$ and
coefficient $-g_{a+1}$.  All earlier profiles are determined by the state,
so both coefficients come from the same recursion parent.  Later splits
retain them.  Both representations carry the identical weight
$f_\eta(X)$, and their divergences cancel.  This also holds for
$a=0,U-1$.

Consequently the summed flow $J$ has divergence $f_\eta/z_-$ on the
minus child, $-f_\eta/z_+$ on the plus child, and zero elsewhere.
The resistance-energy comparison from the start of the proof and
Equation~\eqref{eq:leaf-potential} now eliminate the flow and all
recursive coefficients:
\begin{equation}\label{eq:soft-transport}
 (H_-^\eta-H_+^\eta)^2
 \le\left(\frac1{z_-^\eta}+\frac1{z_+^\eta}
       +\frac{2(U+1)^2\sum_{t>s}c_{st,l}^{\eta,\sigma}}
                    {z_-^\eta z_+^\eta}\right)E_{\eta,\sigma}(H).
\end{equation}

\emph{The limit and conditional repair.}
The extension of $H$ was fixed independently of $\eta$. For this fixed
input all sums are finite.  Since the two limiting
child masses are positive, every quantity in
Equation~\eqref{eq:soft-transport} converges as $\eta\downarrow0$.
We take this limit only after obtaining that inequality; convergence
of the demands, the $g_j$'s, or the flows is unnecessary.
For the limiting $c_{st,l}^\sigma$, Lemma~\ref{lem:repair} injects
state/completion pairs into one of the two child fibers.  Specifically,
since $t$ follows $s$ in row-major order, the receiver is either $s$
itself or lies in a later row, and all earlier assignments are preserved.
The repaired special value is $l$ in the first case and $l-1$ in the
second, with the choice fixed by $s,t$.  Thus
\[
 c_{st,l}^\sigma\le\max(z_-,z_+).
\]
There are fewer than $d$ possible slots $t$.  Multiplication of the
limiting inequality by $\min(z_-,z_+)$ now yields
Equation~\eqref{eq:conditional-transport}.
\end{proof}

\subsection{From conditional transport to variance}

\begin{theorem}[Graph Poincar\'e inequality]\label{thm:poincare}
For every real function $H$ on the feasible small graph,
\begin{equation}\label{eq:graph-poincare}
 \Lambda\Var_{f/\Lambda}H
 \le \left((1+2d^2)4d^2(U+1)^6+2\right)E(H)
 \le 770d^4U^6 E(H).
\end{equation}
In particular, the feasible graph is connected.
\end{theorem}
\begin{proof}
First restrict to transversals.  At a positive-mass context $\sigma$,
put $p_j=z_j/z_\sigma$.  The between-child contribution to its
unnormalized variance is
\begin{equation}\label{eq:child-variance}
 z_\sigma\Var_p(H_j)
 =\frac1{z_\sigma}\sum_{a<b}z_a z_b(H_a-H_b)^2.
\end{equation}
Zero-mass children may be omitted.  For positive endpoints $a<b$,
Lemma~\ref{lem:child-logconcavity} makes every intervening mass at least
$\min(z_a,z_b)$.  Telescoping and Cauchy--Schwarz therefore give
\begin{align*}
 \frac{z_a z_b}{z_\sigma}(H_a-H_b)^2
 &\le \min(z_a,z_b)(b-a)
             \sum_{j=a}^{b-1}(H_j-H_{j+1})^2\\
 &\le U\sum_{j=a}^{b-1}\min(z_j,z_{j+1})(H_j-H_{j+1})^2\\
 &\le U^2 A_0 E_\sigma(H).
\end{align*}
There are at most $(U+1)^2$ pairs.  Thus the contribution in
Equation~\eqref{eq:child-variance} is at most
$(U+1)^4A_0 E_\sigma(H)$.

Iterating the identity
$\Var H=\E(\Var(H\mid j))+\Var(\E(H\mid j))$ down the exposure tree
expresses the full unnormalized transversal variance as the sum of
these contributions.  At any fixed depth, the energies $E_\sigma$
use disjoint edge sets, because distinct contexts assign different
prefix coordinates.  There are at most $d$ internal levels, and
$A_0\le4d(U+1)^2$.  Hence, writing $\Var_{\mathcal T}$ for variance
under the normalized transversal weights,
\begin{equation}\label{eq:transversal-variance}
 Z_0\Var_{\mathcal T}H\le4d^2(U+1)^6E(H).
\end{equation}
If there are no small slots, the transversal space is a singleton and
this inequality holds directly.

Let $\overline H$ be the transversal mean, and compare each defect $X$
to its repair $\mathcal R(X)$.  By Lemma~\ref{lem:repair},
$f(X)\le f(\mathcal R(X))$, the repair map is injective within each
of fewer than $d^2$ defect types, and its edge has conductance $f(X)$.
Distinct defects give distinct repair edges.  It follows that
\begin{align*}
 \Lambda\Var_{f/\Lambda}H
 &\le\sum_{X\in\mathcal X}f(X)(H(X)-\overline H)^2\\
 &\le Z_0\Var_{\mathcal T}H
      +2\sum_{X\notin\mathcal T}f(X)
                         (H(X)-H(\mathcal R(X)))^2\\
 &\hspace{1.5em}+2\sum_{X\notin\mathcal T}f(X)
                         (H(\mathcal R(X))-\overline H)^2\\
 &\le (1+2d^2)Z_0\Var_{\mathcal T}H+2E(H).
\end{align*}
Insert Equation~\eqref{eq:transversal-variance} to obtain the first
bound in Equation~\eqref{eq:graph-poincare}.  Since $U\ge1$,
$1+2d^2\le3d^2$ and $(U+1)^6\le64U^6$ give the second bound.
Finally, the indicator of a nonempty proper connected component would
have positive variance and zero energy, which is impossible.  The
one-state case is already connected.
\end{proof}
